%%% FIELD proof-inclusion
Let \(Q\in\mathcal P_{\mathrm{att}}\).  By Definition
\(\mathrm{def{:}attainable\text{-}law\text{-}class}\), \(Q\) satisfies
consistency, ignorability, bounded outcomes, overlap, propensity H\"older
smoothness, control-regression H\"older smoothness, and CATE H\"older
smoothness, which are all the nondensity member conditions in Definition
\(\mathrm{def{:}full\text{-}law\text{-}class}\).  It remains only to verify
the two density conditions in the latter definition.  Assumption
\(\mathrm{ass{:}uniform\text{-}design}\) gives
\[
 f_Q(x)=1
 \quad\text{for Lebesgue-almost every }x\in\mathcal X.
\tag{1}
\]
The declared ranges \(0<\underline f\le1\) and
\(1\le\overline f<\infty\), together with (1), imply
\[
 f_Q(x)\le\overline f
 \quad\text{for Lebesgue-almost every }x\in\mathcal X
\tag{2}
\]
and
\[
 f_Q(x)\ge\underline f
 \quad\text{for Lebesgue-almost every }
 x\in x_0+[-r_0,r_0]^d.
\tag{3}
\]
Thus \(Q\) also satisfies the global density upper bound and the anchor
local density lower bound.  Every member condition in Definition
\(\mathrm{def{:}full\text{-}law\text{-}class}\) now holds, so
\(Q\in\mathcal P_{\mathrm{full}}\).  Since \(Q\) was arbitrary,
\(\mathcal P_{\mathrm{att}}\subset\mathcal P_{\mathrm{full}}\).

%%% FIELD statement-attainable-bracket
Let \(\mathfrak C_n^{\mathrm{att}}\) be the collection of all replacement
\((\epsilon_n,\delta_n)\)-differentially private random intervals \(C\)
satisfying
\[
 \inf_{Q\in\mathcal P_{\mathrm{att}}}
 \Pr_Q\{\tau_Q^{\mathrm c}(x_0)\in C\}\ge1-\eta,
\]
and define
\[
 \mathcal L_n^*(\mathcal P_{\mathrm{att}})
 =\inf_{C\in\mathfrak C_n^{\mathrm{att}}}
   \sup_{Q\in\mathcal P_{\mathrm{att}}}
   \mathbb E_Q[\operatorname{length}(C)].
\]
If \(L_\pi\ge\kappa\) and \(\delta_n\le c_0/n\), then there are
constants \(c,C>0\) and \(n_0\in\mathbb N\), depending only on the fixed
class constants and \(\eta\), such that, for every \(n\ge n_0\),
\[
 \begin{aligned}
 c\left[n^{-\gamma/(2\gamma+d)}
 \vee\left\{1\wedge(n\epsilon_n)^{-\gamma/(\gamma+d)}\right\}\right]
 &\le \mathcal L_n^*(\mathcal P_{\mathrm{att}})\\
 &\le C\min\left\{
 1\wedge r_{\mathrm{up}},
 1\wedge\left[n^{-\beta/(2\beta+d)}
       \vee N_{\mathrm{priv}}^{-\beta/(\beta+d)}\right]
 \right\}.
 \end{aligned}
\tag{1}
\]
The minimum in (1) can be a strict polynomial improvement over the
higher-order interval alone.  In particular, if \(q\ge\gamma\),
\(\beta<\gamma<2\beta\), and
\(N_{\mathrm{priv}}=n^{c_{\mathrm p}}\), where
\[
 0<c_{\mathrm p}<
 \min\left\{\frac{\gamma+2d}{2\gamma+d},
             \frac{\beta+d}{2\beta+d}\right\},
\tag{2}
\]
then, for all sufficiently large \(n\), the two entries of the upper
minimum are respectively
\[
 n^{-c_{\mathrm p}\gamma/(\gamma+2d)}
 \quad\text{and}\quad
 n^{-c_{\mathrm p}\beta/(\beta+d)},
\]
and their ratio, second over first, converges polynomially to zero.  This
is the best certified bracket here for the attainable class; it makes no
claim that the full-class open frontier is sharp.

%%% FIELD proof-attainable-bracket
We first adapt the two-law argument in Theorem
\(\mathrm{thm{:}full\text{-}class\text{-}ordinary\text{-}lower}\) to
intervals honest only on \(\mathcal P_{\mathrm{att}}\).  Use exactly the
bump \(\varphi\), constants \(p_*=\min\{1/4,L_\mu/2\}\) and \(a_*>0\),
and functions
\[
 g_h(x)=a_*h^\gamma\varphi((x-x_0)/h)
\tag{3}
\]
constructed in that theorem.  Its H\"older calculation gives
\[
 g_h\in\mathcal H^\gamma(L_\tau;\mathcal X),
 \qquad 0\le g_h\le p_*/2,
 \qquad g_h(x_0)=a_*h^\gamma,
\tag{4}
\]
for \(0<h\le r_0\).  Let \(Q_0,Q_h\) be the same full-data laws: under
both, \(X\) is uniform on \(\mathcal X\),
\(A\mid X\sim\operatorname{Bernoulli}(\kappa)\), and \(A\) is
conditionally independent of \((Y(0),Y(1))\) given \(X\); under \(Q_0\)
the two conditionally independent potential outcomes are Bernoulli with
mean \(p_*\), whereas under \(Q_h\) their respective conditional means are
\[
 \mathbb E_{Q_h}[Y(0)\mid X=x]=p_*,
 \qquad
 \mathbb E_{Q_h}[Y(1)\mid X=x]=p_*+g_h(x).
\tag{5}
\]
In both laws set \(Y=AY(1)+(1-A)Y(0)\).

We verify the stronger membership needed here.  The last construction
gives consistency.  The specified conditional independence gives
ignorability.  The potential outcomes have the common support
\(\{0,1\}\), and (4) gives
\[
 0<p_*\le p_*+g_h(x)\le3p_*/2\le3/8<1,
\tag{6}
\]
so the Bernoulli laws are well defined and bounded outcomes holds.  Since
\(0<\kappa<1/2\), the constant propensity satisfies
\[
 \kappa\le\pi_Q(x)=\kappa\le1-\kappa.
\tag{7}
\]
Its H\"older norm is \(\kappa\), so the hypothesis \(L_\pi\ge\kappa\)
gives propensity smoothness.  The control regression is the constant
\(p_*\), whose H\"older norm is at most \(L_\mu/2\), and the contrasts are
zero under \(Q_0\) and \(g_h\) under \(Q_h\), so control-regression and
CATE smoothness follow from the definition of \(p_*\) and (4).  Finally,
the design density equals one almost everywhere, which is exactly
Assumption \(\mathrm{ass{:}uniform\text{-}design}\).  Definition
\(\mathrm{def{:}attainable\text{-}law\text{-}class}\) therefore gives
\[
 Q_0,Q_h\in\mathcal P_{\mathrm{att}}.
\tag{8}
\]

Let \(P_0,P_h\) be their one-record observed laws.  The exact calculations
in Theorem \(\mathrm{thm{:}full\text{-}class\text{-}ordinary\text{-}lower}\)
give constants \(C_\chi,C_v\in(0,\infty)\) such that
\[
 \chi^2(P_h\Vert P_0)=C_\chi h^{2\gamma+d},
 \qquad
 \lVert P_h-P_0\rVert_{\mathrm{TV}}=C_vh^{\gamma+d},
 \qquad
 \left|\tau_{Q_h}^{\mathrm c}(x_0)
       -\tau_{Q_0}^{\mathrm c}(x_0)\right|=a_*h^\gamma.
\tag{9}
\]
That theorem's product-divergence and replacement-private Hamming-transport
steps, using \(\delta_n\le c_0/n\), choose constants
\(c_1,c_2\in(0,r_0]\) for which
\[
 h_1=c_1n^{-1/(2\gamma+d)},
 \qquad
 h_2=c_2\{1\wedge(n\epsilon_n)^{-1}\}^{1/(\gamma+d)}
\tag{10}
\]
have the following property: for every replacement
\((\epsilon_n,\delta_n)\)-private release kernel \(\mathsf K\),
\[
 \lVert P_{h_j}^n\mathsf K-P_0^n\mathsf K\rVert_{\mathrm{TV}}
 \le1/8,
 \qquad j\in\{1,2\}.
\tag{11}
\]
No full-class honesty premise enters (9)--(11).

Now take any \(C\in\mathfrak C_n^{\mathrm{att}}\), and let \(\mathsf K\)
be its release kernel.  For either \(h=h_1\) or \(h=h_2\), write
\(\theta_0=\tau_{Q_0}^{\mathrm c}(x_0)\),
\(\theta_h=\tau_{Q_h}^{\mathrm c}(x_0)\),
\(\nu_0=P_0^n\mathsf K\), and \(\nu_h=P_h^n\mathsf K\).  By (8) and
attainable-class honesty,
\[
 \nu_0\{\theta_0\in C\}\ge1-\eta,
 \qquad
 \nu_h\{\theta_h\in C\}\ge1-\eta.
\tag{12}
\]
Combining (11)--(12) and the union bound gives
\[
 \begin{aligned}
 \nu_0\{\theta_0\in C,\theta_h\in C\}
 &\ge (1-\eta)+(1-\eta-1/8)-1\\
 &=1-2\eta-1/8>3/8,
 \end{aligned}
\tag{13}
\]
where the final inequality uses \(0<\eta<1/4\).  On the event in (13),
an interval containing both target values has length at least
\(a_*h^\gamma\) by (9).  Hence
\[
 \mathbb E_{Q_0}[\operatorname{length}(C)]
 \ge(3a_*/8)h^\gamma.
\tag{14}
\]
Applying (14) at both values in (10), taking their maximum, then taking
the supremum over \(Q\in\mathcal P_{\mathrm{att}}\) and the infimum over
\(C\in\mathfrak C_n^{\mathrm{att}}\), proves the lower inequality in (1).

For the upper inequality, put
\[
 a_n=1\wedge r_{\mathrm{up}},
 \qquad
 b_n=1\wedge\left[n^{-\beta/(2\beta+d)}
          \vee N_{\mathrm{priv}}^{-\beta/(\beta+d)}\right].
\tag{15}
\]
Choose the deterministic minimizing \((h,k)\) from Theorem
\(\mathrm{thm{:}phase\text{-}sensitive\text{-}expected\text{-}length}\)
and any fixed deterministic pilot dimensions.  Theorem
\(\mathrm{thm{:}honest\text{-}private\text{-}pointwise\text{-}interval\text{-}unrestricted\text{-}rank}\)
puts the resulting higher-order interval \(C_n^{\mathrm{HO}}\) in
\(\mathfrak C_n^{\mathrm{att}}\), while the phase-sensitive theorem gives
\[
 \sup_{Q\in\mathcal P_{\mathrm{att}}}
 \mathbb E_Q[\operatorname{length}(C_n^{\mathrm{HO}})]
 \le C_{\mathrm{HO}}a_n.
\tag{16}
\]
The direct interval \(C_n^{\mathrm{dir}}(h_n)\) in Theorem
\(\mathrm{thm{:}full\text{-}class\text{-}direct\text{-}regression\text{-}upper}\)
is replacement private and honest on \(\mathcal P_{\mathrm{full}}\).
Proposition \(\mathrm{prop{:}attainable\text{-}inclusion}\) implies
\[
 \inf_{Q\in\mathcal P_{\mathrm{att}}}
 \Pr_Q\{\tau_Q^{\mathrm c}(x_0)\in C_n^{\mathrm{dir}}(h_n)\}
 \ge1-\eta
\tag{17}
\]
and
\[
 \sup_{Q\in\mathcal P_{\mathrm{att}}}
 \mathbb E_Q[\operatorname{length}\{C_n^{\mathrm{dir}}(h_n)\}]
 \le C_{\mathrm{dir}}b_n.
\tag{18}
\]
Thus both candidates belong to \(\mathfrak C_n^{\mathrm{att}}\).  Equivalently,
one may use the public deterministic selector which releases only
\(C_n^{\mathrm{HO}}\) when \(a_n\le b_n\), and releases only
\(C_n^{\mathrm{dir}}(h_n)\) otherwise.  The selector depends only on public
parameters, runs one mechanism, and hence incurs no privacy composition.
Equations (16)--(18) yield
\[
 \mathcal L_n^*(\mathcal P_{\mathrm{att}})
 \le \min\{C_{\mathrm{HO}}a_n,C_{\mathrm{dir}}b_n\}
 \le \max\{C_{\mathrm{HO}},C_{\mathrm{dir}}\}\min\{a_n,b_n\},
\tag{19}
\]
which is the upper inequality in (1).

It remains to verify the strict-improvement regime.  Suppose the conditions
following (1) hold.  Since \(q\ge\gamma\), one has \(s=q/2\ge\gamma/2\),
and
\[
 \frac{\gamma}{2}>
 \frac{d/4}{1+d/(2\gamma)}.
\tag{20}
\]
Definition \(\mathrm{def{:}phase\text{-}rate}\) and (20) therefore give
\[
 r_{\mathrm{np}}=n^{-\gamma/(2\gamma+d)},
 \quad
 r_{\mathrm{reg}}=n^{-c_{\mathrm p}\gamma/(\gamma+d)},
 \quad
 r_{\mathrm{ho}}=n^{-c_{\mathrm p}\gamma/(\gamma+2d)}.
\tag{21}
\]
The first restriction in (2) makes the last term in (21) dominate
\(r_{\mathrm{np}}\), and it also dominates \(r_{\mathrm{reg}}\).  Hence,
for \(n>1\),
\[
 a_n=n^{-c_{\mathrm p}\gamma/(\gamma+2d)}.
\tag{22}
\]
The second restriction in (2) makes the privacy term dominate the sampling
term in \(b_n\), so
\[
 b_n=n^{-c_{\mathrm p}\beta/(\beta+d)}.
\tag{23}
\]
Finally, \(\gamma<2\beta\) is equivalent to
\[
 \frac{\beta}{\beta+d}>\frac{\gamma}{\gamma+2d}.
\tag{24}
\]
Equations (22)--(24) imply
\[
 \frac{b_n}{a_n}
 =n^{-c_{\mathrm p}\{\beta/(\beta+d)-\gamma/(\gamma+2d)\}}
 \longrightarrow0,
\]
proving the stated strict polynomial improvement.  All lower and upper
comparisons above are on \(\mathcal P_{\mathrm{att}}\); nothing in the
argument closes the separate full-class open frontier.
